% Options for packages loaded elsewhere
% Options for packages loaded elsewhere
\PassOptionsToPackage{unicode}{hyperref}
\PassOptionsToPackage{hyphens}{url}
\PassOptionsToPackage{dvipsnames,svgnames,x11names}{xcolor}
%
\documentclass[
  11pt,
  a4paper,
]{article}
\usepackage{xcolor}
\usepackage[margin=2.5cm]{geometry}
\usepackage{amsmath,amssymb}
\setcounter{secnumdepth}{5}
\usepackage{iftex}
\ifPDFTeX
  \usepackage[T1]{fontenc}
  \usepackage[utf8]{inputenc}
  \usepackage{textcomp} % provide euro and other symbols
\else % if luatex or xetex
  \usepackage{unicode-math} % this also loads fontspec
  \defaultfontfeatures{Scale=MatchLowercase}
  \defaultfontfeatures[\rmfamily]{Ligatures=TeX,Scale=1}
\fi
\usepackage{lmodern}
\ifPDFTeX\else
  % xetex/luatex font selection
\fi
% Use upquote if available, for straight quotes in verbatim environments
\IfFileExists{upquote.sty}{\usepackage{upquote}}{}
\IfFileExists{microtype.sty}{% use microtype if available
  \usepackage[]{microtype}
  \UseMicrotypeSet[protrusion]{basicmath} % disable protrusion for tt fonts
}{}
\makeatletter
\@ifundefined{KOMAClassName}{% if non-KOMA class
  \IfFileExists{parskip.sty}{%
    \usepackage{parskip}
  }{% else
    \setlength{\parindent}{0pt}
    \setlength{\parskip}{6pt plus 2pt minus 1pt}}
}{% if KOMA class
  \KOMAoptions{parskip=half}}
\makeatother
% Make \paragraph and \subparagraph free-standing
\makeatletter
\ifx\paragraph\undefined\else
  \let\oldparagraph\paragraph
  \renewcommand{\paragraph}{
    \@ifstar
      \xxxParagraphStar
      \xxxParagraphNoStar
  }
  \newcommand{\xxxParagraphStar}[1]{\oldparagraph*{#1}\mbox{}}
  \newcommand{\xxxParagraphNoStar}[1]{\oldparagraph{#1}\mbox{}}
\fi
\ifx\subparagraph\undefined\else
  \let\oldsubparagraph\subparagraph
  \renewcommand{\subparagraph}{
    \@ifstar
      \xxxSubParagraphStar
      \xxxSubParagraphNoStar
  }
  \newcommand{\xxxSubParagraphStar}[1]{\oldsubparagraph*{#1}\mbox{}}
  \newcommand{\xxxSubParagraphNoStar}[1]{\oldsubparagraph{#1}\mbox{}}
\fi
\makeatother


\usepackage{longtable,booktabs,array}
\newcounter{none} % for unnumbered tables
\usepackage{calc} % for calculating minipage widths
% Correct order of tables after \paragraph or \subparagraph
\usepackage{etoolbox}
\makeatletter
\patchcmd\longtable{\par}{\if@noskipsec\mbox{}\fi\par}{}{}
\makeatother
% Allow footnotes in longtable head/foot
\IfFileExists{footnotehyper.sty}{\usepackage{footnotehyper}}{\usepackage{footnote}}
\makesavenoteenv{longtable}
\usepackage{graphicx}
\makeatletter
\newsavebox\pandoc@box
\newcommand*\pandocbounded[1]{% scales image to fit in text height/width
  \sbox\pandoc@box{#1}%
  \Gscale@div\@tempa{\textheight}{\dimexpr\ht\pandoc@box+\dp\pandoc@box\relax}%
  \Gscale@div\@tempb{\linewidth}{\wd\pandoc@box}%
  \ifdim\@tempb\p@<\@tempa\p@\let\@tempa\@tempb\fi% select the smaller of both
  \ifdim\@tempa\p@<\p@\scalebox{\@tempa}{\usebox\pandoc@box}%
  \else\usebox{\pandoc@box}%
  \fi%
}
% Set default figure placement to htbp
\def\fps@figure{htbp}
\makeatother





\setlength{\emergencystretch}{3em} % prevent overfull lines

\providecommand{\tightlist}{%
  \setlength{\itemsep}{0pt}\setlength{\parskip}{0pt}}



 


\usepackage{booktabs}
\usepackage{float}
\usepackage{caption}
\captionsetup{font=small,labelfont=bf}
\usepackage{fancyhdr}
\usepackage{xcolor}
\pagestyle{fancy}
\fancyhead[L]{\small BTC-QUANT v4.4-v4.5}
\fancyhead[R]{\small Performance Logbook}
\fancyfoot[C]{\thepage}
\renewcommand{\headrulewidth}{0.4pt}
\definecolor{win}{RGB}{34, 139, 34}
\definecolor{loss}{RGB}{220, 20, 60}
\makeatletter
\@ifpackageloaded{caption}{}{\usepackage{caption}}
\AtBeginDocument{%
\ifdefined\contentsname
  \renewcommand*\contentsname{Table of contents}
\else
  \newcommand\contentsname{Table of contents}
\fi
\ifdefined\listfigurename
  \renewcommand*\listfigurename{List of Figures}
\else
  \newcommand\listfigurename{List of Figures}
\fi
\ifdefined\listtablename
  \renewcommand*\listtablename{List of Tables}
\else
  \newcommand\listtablename{List of Tables}
\fi
\ifdefined\figurename
  \renewcommand*\figurename{Figure}
\else
  \newcommand\figurename{Figure}
\fi
\ifdefined\tablename
  \renewcommand*\tablename{Table}
\else
  \newcommand\tablename{Table}
\fi
}
\@ifpackageloaded{float}{}{\usepackage{float}}
\floatstyle{ruled}
\@ifundefined{c@chapter}{\newfloat{codelisting}{h}{lop}}{\newfloat{codelisting}{h}{lop}[chapter]}
\floatname{codelisting}{Listing}
\newcommand*\listoflistings{\listof{codelisting}{List of Listings}}
\makeatother
\makeatletter
\makeatother
\makeatletter
\@ifpackageloaded{caption}{}{\usepackage{caption}}
\@ifpackageloaded{subcaption}{}{\usepackage{subcaption}}
\makeatother
\usepackage{bookmark}
\IfFileExists{xurl.sty}{\usepackage{xurl}}{} % add URL line breaks if available
\urlstyle{same}
\hypersetup{
  pdftitle={BTC-QUANT Performance Logbook},
  pdfauthor={BTC-QUANT Execution Layer},
  pdfkeywords={BTC-QUANT, live trading, HestonStrategy, performance
analysis, trading bot},
  colorlinks=true,
  linkcolor={blue},
  filecolor={Maroon},
  citecolor={Blue},
  urlcolor={Blue},
  pdfcreator={LaTeX via pandoc}}


\title{BTC-QUANT Performance Logbook}
\usepackage{etoolbox}
\makeatletter
\providecommand{\subtitle}[1]{% add subtitle to \maketitle
  \apptocmd{\@title}{\par {\large #1 \par}}{}{}
}
\makeatother
\subtitle{Live Trading Analysis: March - April 2026}
\author{BTC-QUANT Execution Layer}
\date{2026-04-17}
\begin{document}
\maketitle
\begin{abstract}
Laporan ini mendokumentasikan performa live trading BTC-QUANT dari 10
Maret hingga 14 April 2026 menggunakan strategi HestonStrategy dengan
parameter margin \$20 dan leverage 7x. Analisis mencakup 54 posisi
dengan win rate 75.9\% dan total PnL \$20.37 (adjusted). Laporan ini
secara khusus menonjolkan lesson learned dari intervensi manual yang
merugikan, dengan kesimpulan utama: percaya pada edge sistem dan hindari
intervensi manual.
\end{abstract}

\renewcommand*\contentsname{Table of contents}
{
\setcounter{tocdepth}{2}
\tableofcontents
}

\section{Glosarium \& Konsep Utama}\label{glosarium-konsep-utama}

{\def\LTcaptype{none} % do not increment counter
\begin{longtable}[]{@{}
  >{\raggedright\arraybackslash}p{(\linewidth - 2\tabcolsep) * \real{0.4286}}
  >{\raggedright\arraybackslash}p{(\linewidth - 2\tabcolsep) * \real{0.5714}}@{}}
\toprule\noalign{}
\begin{minipage}[b]{\linewidth}\raggedright
Istilah
\end{minipage} & \begin{minipage}[b]{\linewidth}\raggedright
Keterangan
\end{minipage} \\
\midrule\noalign{}
\endhead
\bottomrule\noalign{}
\endlastfoot
\textbf{PnL (Profit \& Loss)} & Keuntungan atau kerugian bersih dari
posisi yang ditutup. \\
\textbf{Win Rate (WR)} & Persentase posisi profitable dari total
posisi. \\
\textbf{Intervensi Manual} & Tindakan trader menggeser SL/TP atau
menutup posisi secara manual, mengesampingkan sistem otomatis. \\
\textbf{Signal Delay} & Kondisi di mana posisi terbuka tapi tidak
tereksekusi TP/SL dalam waktu yang diharapkan. \\
\textbf{Time-based Exit} & Mekanisme auto-close posisi setelah durasi
tertentu tanpa trigger TP/SL. \\
\textbf{HestonStrategy} & Strategi dengan SL/TP adaptif berbasis ATR
dari model volatilitas Heston. \\
\textbf{Edge} & Keunggulan statistik sistem trading yang teruji secara
backtest. \\
\end{longtable}
}

\section{Ringkasan Eksekutif (per Dashboard
Lighter)}\label{ringkasan-eksekutif-per-dashboard-lighter}

\begin{longtable}[]{@{}ll@{}}
\caption{Ringkasan metrik performa dari dashboard
Lighter.}\label{tbl-executive}\tabularnewline
\toprule\noalign{}
Metrik & Nilai \\
\midrule\noalign{}
\endfirsthead
\toprule\noalign{}
Metrik & Nilai \\
\midrule\noalign{}
\endhead
\bottomrule\noalign{}
\endlastfoot
\textbf{PnL} & \textbf{\$14.11} \\
\textbf{Volume} & \$25,972.88 \\
\textbf{Return Percentage} & 14.19\% \\
\textbf{Average Daily PnL} & \$0.35 \\
\textbf{PnL Volatility} & \$2.15 \\
\textbf{Sharpe Ratio} & 3.13 \\
\textbf{Maximum Drawdown} & \$6.53 \\
\textbf{Estimated Fees Saved} & \$12.98 \\
\end{longtable}

\textbf{Note:} Dashboard menunjukkan PnL \$14.11. Setelah adjustment
untuk bug 18 Mar (loss -\$6.30 diabaikan) dan filter 15 Mar (uji coba),
\textbf{PnL aktual adalah \$20.37} dengan \textbf{54 posisi} dan
\textbf{win rate 75.9\%}.

\textbf{Expected Value (Potensi):} \textasciitilde\$35 (tanpa intervensi
manual dan sizing error)

\newpage

\section{Adjustments \& Data Filter}\label{adjustments-data-filter}

\subsection{Data yang Diabaikan}\label{data-yang-diabaikan}

\begin{longtable}[]{@{}
  >{\raggedright\arraybackslash}p{(\linewidth - 6\tabcolsep) * \real{0.3000}}
  >{\raggedright\arraybackslash}p{(\linewidth - 6\tabcolsep) * \real{0.2667}}
  >{\raggedright\arraybackslash}p{(\linewidth - 6\tabcolsep) * \real{0.1667}}
  >{\raggedright\arraybackslash}p{(\linewidth - 6\tabcolsep) * \real{0.2667}}@{}}
\caption{Data yang diexclude dari perhitungan
final.}\label{tbl-excluded}\tabularnewline
\toprule\noalign{}
\begin{minipage}[b]{\linewidth}\raggedright
Tanggal
\end{minipage} & \begin{minipage}[b]{\linewidth}\raggedright
Posisi
\end{minipage} & \begin{minipage}[b]{\linewidth}\raggedright
PnL
\end{minipage} & \begin{minipage}[b]{\linewidth}\raggedright
Alasan
\end{minipage} \\
\midrule\noalign{}
\endfirsthead
\toprule\noalign{}
\begin{minipage}[b]{\linewidth}\raggedright
Tanggal
\end{minipage} & \begin{minipage}[b]{\linewidth}\raggedright
Posisi
\end{minipage} & \begin{minipage}[b]{\linewidth}\raggedright
PnL
\end{minipage} & \begin{minipage}[b]{\linewidth}\raggedright
Alasan
\end{minipage} \\
\midrule\noalign{}
\endhead
\bottomrule\noalign{}
\endlastfoot
\textbf{15 Mar 2026} & 5 & -\$0.04 & Uji coba otomasi bot, tidak
mencerminkan strategi actual. \\
\textbf{18 Mar 2026} & 1 & -\$6.30 & Bug autentikasi SL order, dihitung
sebagai loss tapi PnL = 0. \\
\end{longtable}

\textbf{Rationale:} 15 Mar merupakan testing phase dengan intervensi
manual intensif. 18 Mar adalah loss akibat bug autentikasi SL pada
Lighter gateway (sudah diperbaiki di commit \texttt{c3107f9}), bukan
kesalahan strategi.

\section{Analisis Per Periode}\label{analisis-per-periode}

\subsection{Maret 2026: Fase Setup \&
Debugging}\label{maret-2026-fase-setup-debugging}

\begin{longtable}[]{@{}lllll@{}}
\caption{Detail performa harian Maret
2026.}\label{tbl-march}\tabularnewline
\toprule\noalign{}
Tanggal & Posisi & PnL & Status & Catatan \\
\midrule\noalign{}
\endfirsthead
\toprule\noalign{}
Tanggal & Posisi & PnL & Status & Catatan \\
\midrule\noalign{}
\endhead
\bottomrule\noalign{}
\endlastfoot
10 Mar & 1 & -\$0.12 & ✅ & Fase manual awal \\
11 Mar & 2 & +\$0.86 & ✅ & -- \\
12 Mar & 3 & -\$0.91 & ✅ & -- \\
13 Mar & 5 & +\$0.83 & ✅ & -- \\
16 Mar & 2 & +\$1.65 & ✅ & Start fase otomatis \\
17 Mar & 7 & +\$1.30 & ✅ & -- \\
22 Mar & 3 & +\$0.28 & ✅ & -- \\
23 Mar & 1 & +\$0.26 & ✅ & -- \\
26 Mar & 2 & +\$0.42 & ✅ & -- \\
27 Mar & 1 & +\$0.26 & ✅ & -- \\
28 Mar & 1 & +\$0.11 & ✅ & -- \\
29 Mar & 4 & \textbf{+\$3.14} & ✅ & Strong short \\
31 Mar & 1 & -\$1.91 & ✅ & -- \\
\end{longtable}

\textbf{Total Maret:} +\$6.17 (33 posisi, adjusted)

\begin{center}\rule{0.5\linewidth}{0.5pt}\end{center}

\subsection{April 2026: Fase
Operasional}\label{april-2026-fase-operasional}

\begin{longtable}[]{@{}lllll@{}}
\caption{Detail performa harian April
2026.}\label{tbl-april}\tabularnewline
\toprule\noalign{}
Tanggal & Posisi & PnL & Status & Catatan \\
\midrule\noalign{}
\endfirsthead
\toprule\noalign{}
Tanggal & Posisi & PnL & Status & Catatan \\
\midrule\noalign{}
\endhead
\bottomrule\noalign{}
\endlastfoot
2 Apr & 1 & +\$3.54 & ✅ & -- \\
\textbf{3 Apr} & 2 & +\$6.35 & ⚠️ & Interupsi manual, profit \textless{}
target \\
\textbf{4 Apr} & 2 & +\$0.18 & ❌ & \textbf{Interupsi manual → profit
minim} \\
5 Apr & 1 & -\$2.99 & ✅ & Loss normal \\
6 Apr & 2 & -\$3.52 & ✅ & -- \\
\textbf{7 Apr} & 1 & +\$0.12 & ⚠️ & Interupsi manual, profit minim \\
8 Apr & 1 & +\$3.36 & ✅ & -- \\
9 Apr & 2 & +\$3.81 & ✅ & -- \\
\textbf{10 Apr} & 2 & +\$4.00 & ⚠️ & Interupsi manual, profit
\textless{} target \\
11 Apr & 2 & -\$2.93 & ✅ & -- \\
\textbf{12 Apr} & 1 & -\$2.52 & ⚠️ & SL manual adjustment, sizing
kecil \\
13 Apr & 2 & +\$3.31 & ✅ & Sizing kecil \\
14 Apr & 1 & +\$1.50 & ✅ & Sizing kecil \\
\end{longtable}

\textbf{Total April:} +\$14.21 (20 posisi)

\begin{center}\rule{0.5\linewidth}{0.5pt}\end{center}

\section{Analisis Intervensi Manual}\label{analisis-intervensi-manual}

\subsection{Konteks Intervensi}\label{konteks-intervensi}

\textbf{Rationale Intervensi:} \textgreater{} Ketika signal sudah lebih
dari 8 jam tapi belum close posisi (TP/SL), maka SL digeser lebih
sedikit untuk mengurangi loss karena signal dinilai tidak lagi relevan
di timeframe tersebut.

\subsection{Impact Interupsi Signal
Delay}\label{impact-interupsi-signal-delay}

\begin{longtable}[]{@{}
  >{\raggedright\arraybackslash}p{(\linewidth - 8\tabcolsep) * \real{0.1525}}
  >{\raggedright\arraybackslash}p{(\linewidth - 8\tabcolsep) * \real{0.2034}}
  >{\raggedright\arraybackslash}p{(\linewidth - 8\tabcolsep) * \real{0.2203}}
  >{\raggedright\arraybackslash}p{(\linewidth - 8\tabcolsep) * \real{0.3051}}
  >{\raggedright\arraybackslash}p{(\linewidth - 8\tabcolsep) * \real{0.1186}}@{}}
\caption{Impact financial dari intervensi
manual.}\label{tbl-intervention}\tabularnewline
\toprule\noalign{}
\begin{minipage}[b]{\linewidth}\raggedright
Tanggal
\end{minipage} & \begin{minipage}[b]{\linewidth}\raggedright
PnL Aktual
\end{minipage} & \begin{minipage}[b]{\linewidth}\raggedright
Estimasi TP
\end{minipage} & \begin{minipage}[b]{\linewidth}\raggedright
Opportunity Cost
\end{minipage} & \begin{minipage}[b]{\linewidth}\raggedright
Hasil
\end{minipage} \\
\midrule\noalign{}
\endfirsthead
\toprule\noalign{}
\begin{minipage}[b]{\linewidth}\raggedright
Tanggal
\end{minipage} & \begin{minipage}[b]{\linewidth}\raggedright
PnL Aktual
\end{minipage} & \begin{minipage}[b]{\linewidth}\raggedright
Estimasi TP
\end{minipage} & \begin{minipage}[b]{\linewidth}\raggedright
Opportunity Cost
\end{minipage} & \begin{minipage}[b]{\linewidth}\raggedright
Hasil
\end{minipage} \\
\midrule\noalign{}
\endhead
\bottomrule\noalign{}
\endlastfoot
3 Apr & +\$6.35 & \textasciitilde\$8.68 & -\$2.33 & Profit, tapi
\textless{} target \\
\textbf{4 Apr} & \textbf{+\$0.18} & \textasciitilde\$8.68 &
\textbf{-\$8.50} & ❌ \textbf{Harusnya 2 TP, malah profit minim} \\
7 Apr & +\$0.12 & \textasciitilde\$4.34 & -\$4.22 & Profit minim \\
10 Apr & +\$4.00 & \textasciitilde\$8.68 & -\$4.68 & Profit, tapi
\textless{} target \\
\end{longtable}

\textbf{Total Opportunity Cost:} - Aktual: +\$10.65 - Potensi:
\textasciitilde\$25.38 - \textbf{Profit yang hilang:
\textasciitilde\$15}

\begin{center}\rule{0.5\linewidth}{0.5pt}\end{center}

\newpage

\subsection{Sizing Error (12-14 April)}\label{sizing-error-12-14-april}

\begin{longtable}[]{@{}lllll@{}}
\caption{Impact sizing error pada
performa.}\label{tbl-sizing}\tabularnewline
\toprule\noalign{}
Tanggal & PnL & Size Aktual & Size Target & Impact \\
\midrule\noalign{}
\endfirsthead
\toprule\noalign{}
Tanggal & PnL & Size Aktual & Size Target & Impact \\
\midrule\noalign{}
\endhead
\bottomrule\noalign{}
\endlastfoot
12 Apr & -\$2.52 & \$147 & \$500 & Loss lebih kecil karena sizing \\
13 Apr & +\$3.31 & \$151 & \$500 & Profit 30\% dari potensi \\
14 Apr & +\$1.50 & \$151 & \$500 & Profit 30\% dari potensi \\
\end{longtable}

\textbf{Note:} 12 Apr SL manual adjustment berhasil menyelamatkan dari
loss lebih besar saat harga anjlok. Namun, sizing kecil mengurangkan
profit potential.

\begin{center}\rule{0.5\linewidth}{0.5pt}\end{center}

\section{Performa by Side (Long vs
Short)}\label{performa-by-side-long-vs-short}

\begin{longtable}[]{@{}llll@{}}
\caption{Perbandingan performa Long vs
Short.}\label{tbl-byside}\tabularnewline
\toprule\noalign{}
Arah & Win/Loss & Win Rate & PnL \\
\midrule\noalign{}
\endfirsthead
\toprule\noalign{}
Arah & Win/Loss & Win Rate & PnL \\
\midrule\noalign{}
\endhead
\bottomrule\noalign{}
\endlastfoot
\textbf{Long} & 27W / 9L & 75.0\% & +\$7.82 \\
\textbf{Short} & 14W / 3L & \textbf{82.4\%} & \textbf{+\$12.55} \\
\end{longtable}

\textbf{Insight:} Short strategy jauh lebih superior dengan win rate
82.4\% dan PnL \$12.55 vs \$7.82. Ini menunjukkan edge yang lebih kuat
pada kondisi bearish/trend down.

\begin{center}\rule{0.5\linewidth}{0.5pt}\end{center}

\section{Lessons Learned}\label{lessons-learned}

\subsection{Pelajaran \#1: Percaya pada
Edge}\label{pelajaran-1-percaya-pada-edge}

\begin{quote}
\textbf{``Percaya pada edge dan jangan lakukan intervensi manual''}
\end{quote}

\textbf{Evidence:} - \textbf{4 Apr:} Intervensi → Profit +\$0.18
(harusnya 2 TP \textasciitilde\$8.68) - \textbf{3, 7, 10 Apr:}
Intervensi → Profit \textless{} target - \textbf{Cost of intervention:
\textasciitilde\$15}

\begin{center}\rule{0.5\linewidth}{0.5pt}\end{center}

\subsection{Pelajaran \#2: Consistent
Sizing}\label{pelajaran-2-consistent-sizing}

\textbf{Issue:} 12-14 Apr sizing hanya \textasciitilde30\% dari target
(\$150 vs \$500 notional)

\textbf{Impact:} - Profit dikurangi \textasciitilde70\% - Loss juga
dikurangi \textasciitilde70\% (12 Apr SL manual)

\textbf{Action:} Fix position sizing logic di bot configuration.

\begin{center}\rule{0.5\linewidth}{0.5pt}\end{center}

\subsection{Pelajaran \#3: SL
Management}\label{pelajaran-3-sl-management}

\textbf{Positif:} 12 Apr SL manual adjustment berhasil menyelamatkan
dari loss lebih besar.

\textbf{Negatif:} Intervensi manual justru merusak edge pada 4 Apr.

\textbf{Rule:} - Jika signal valid dan SL belum hit → biarkan - Jika
signal expired (\textgreater8 jam) → gunakan \textbf{time-exit}, bukan
SL adjustment

\section{Riwayat Perbaikan (Setelah 27
Maret)}\label{riwayat-perbaikan-setelah-27-maret}

\subsection{29 Maret 2026 --- Critical
Fixes}\label{maret-2026-critical-fixes}

{\def\LTcaptype{none} % do not increment counter
\begin{longtable}[]{@{}llll@{}}
\toprule\noalign{}
ID & Deskripsi & Severity & Commit \\
\midrule\noalign{}
\endhead
\bottomrule\noalign{}
\endlastfoot
\textbf{FIX-13} & Re-entry Guard & \textbf{CRITICAL} &
\texttt{ddc6309} \\
\textbf{FIX-14} & Koreksi Database dari CSV & MEDIUM & --- \\
\textbf{FIX-15} & Switch ke FixedStrategy & MEDIUM & \texttt{663ba55} \\
\textbf{FIX-16} & SL/TP Limit Price Slippage Buffer & HIGH &
\texttt{7330b4c} \\
\textbf{FIX-17} & filled\_price dari SDK Response & MEDIUM &
\texttt{7330b4c} \\
\end{longtable}
}

\subsubsection{FIX-13: Re-entry Guard
(Critical)}\label{fix-13-re-entry-guard-critical}

\textbf{Masalah:} Saat posisi baru saja ditutup,
\texttt{sync\_position\_status()} langsung diikuti
\texttt{process\_signal()} di cycle yang sama, menyebabkan bot langsung
membuka posisi baru (double/triple entry).

\textbf{Solusi:} - Return \texttt{True} hanya jika posisi baru saja
ditutup di cycle ini - Caller skip \texttt{process\_signal()} jika
return \texttt{True}

\textbf{Impact:} Menghilangkan bug double/triple entry.

\subsubsection{FIX-14: Koreksi Database dari
CSV}\label{fix-14-koreksi-database-dari-csv}

\textbf{Masalah:} Data DuckDB tidak akurat --- berisi 10 record campuran
dari bug double entry dan data tidak lengkap.

\textbf{Solusi:} Rekonstruksi database dari CSV export Lighter (ground
truth): - Hapus 10 record lama - Insert ulang 9 trade valid berdasarkan
CSV

\textbf{Result:} Total PnL direkonstruksi: \textbf{+\$2.18 USDT}

\subsubsection{FIX-15: Switch ke
FixedStrategy}\label{fix-15-switch-ke-fixedstrategy}

\textbf{Masalah:} HestonStrategy menghasilkan TP terlalu jauh (ATR × 2.1
≈ 3.1\%), posisi tidak hit TP dan harus ditutup manual.

\textbf{Solusi:} Switch ke FixedStrategy (Golden v4.4) dengan parameter:
- SL\_PCT: 1.333\% - TP\_PCT: 0.71\% - LEVERAGE: 7x - MARGIN\_USD: \$20

\subsubsection{FIX-16: SL/TP Limit Price Slippage
Buffer}\label{fix-16-sltp-limit-price-slippage-buffer}

\textbf{Masalah:} Limit price SL/TP di-set sama dengan trigger price,
order tidak fill saat market gap/slippage.

\textbf{Solusi:} Tambah slippage buffer: - SL order: buffer 0.5\% - TP
order: buffer 0.3\%

\subsubsection{FIX-17: filled\_price dari SDK
Response}\label{fix-17-filled_price-dari-sdk-response}

\textbf{Masalah:} \texttt{filled\_price} di-set dari intended price,
bukan actual fill price dari SDK.

\textbf{Solusi:} Baca \texttt{avg\_execution\_price} dari SDK response
object.

\section{Rekomendasi Improvement}\label{rekomendasi-improvement}

\subsection{Immediate (High Priority)}\label{immediate-high-priority}

\begin{enumerate}
\def\labelenumi{\arabic{enumi}.}
\tightlist
\item
  \textbf{Time-based Exit:} Auto-close posisi setelah 8-12 jam tanpa
  TP/SL
\item
  \textbf{Fix Sizing:} Pastikan margin \& leverage sesuai config
  (\$20×7x atau \$100×5x)
\item
  \textbf{Remove Manual Override:} Disable ability untuk manual close
  dari dashboard
\end{enumerate}

\subsection{Medium Term}\label{medium-term}

\begin{enumerate}
\def\labelenumi{\arabic{enumi}.}
\tightlist
\item
  \textbf{Signal Decay:} Reduce conviction score setiap candle tanpa
  trigger
\item
  \textbf{Trailing SL:} Implement break-even SL setelah 50\% TP tercapai
\item
  \textbf{Market Regime Detection:} Reduce size atau pause trading saat
  choppy/volatile
\end{enumerate}

\subsection{Long Term}\label{long-term}

\begin{enumerate}
\def\labelenumi{\arabic{enumi}.}
\tightlist
\item
  \textbf{Walk-forward Validation:} Retest strategi dengan intervensi
  manual di-backtest
\item
  \textbf{Feature Analysis:} Apa yang membuat short lebih profitable?
\item
  \textbf{Correlation Analysis:} L1-L3 alignment untuk improve entry
  quality
\end{enumerate}

\begin{center}\rule{0.5\linewidth}{0.5pt}\end{center}

\section{Kesimpulan}\label{kesimpulan}

\begin{enumerate}
\def\labelenumi{\arabic{enumi}.}
\tightlist
\item
  \textbf{Trust the System:} Intervensi manual merugikan dengan
  opportunity cost \textasciitilde\$15.
\item
  \textbf{Short Superiority:} Win rate 82.4\% vs 75\% pada short
  menunjukkan edge lebih kuat.
\item
  \textbf{Sizing Matters:} Inconsistent sizing mengurangi profit
  potential signifikan.
\item
  \textbf{Time-Exit \textgreater{} SL Adjustment:} Mekanisme auto-close
  berdasarkan waktu lebih baik daripada manual SL adjustment.
\end{enumerate}

\textbf{Final PnL Potential:} \$20.37 (aktual) + \$15 (opportunity) =
\textbf{\textasciitilde\$35}

\begin{center}\rule{0.5\linewidth}{0.5pt}\end{center}

\emph{Dokumen ini dibuat untuk mendokumentasikan lesson learned dan
mencegah repeat of mistakes. Terakhir diupdate: 17 April 2026.}




\end{document}
